\documentclass[10pt]{article}

\usepackage[margin=0.9in]{geometry}
\usepackage[T1]{fontenc}
\usepackage{charter}
\usepackage{microtype}
\usepackage{parskip}
\usepackage{amsmath,amssymb}
\usepackage[hidelinks]{hyperref}

\newcommand{\Q}{\mathbb{Q}}
\newcommand{\repourl}{\url{https://github.com/mmjerge/galois-ml}}

\title{\Large\bfseries Evaluating Agents on Scientific Discovery\\
\large A Benchmark with Provable Ground Truth from Arithmetic Statistics}
\author{Michael Jerge}
\date{August 2026 --- working draft}

\begin{document}
\maketitle
\thispagestyle{empty}

\subsection*{Motivation}

Evaluation of LLM-based agents on scientific discovery is
methodologically circular: scoring a discovered law presupposes
knowledge of the true law, which is unavailable precisely for the
problems of interest. Current discovery benchmarks substitute LLM
judges or rubric matching, inheriting the reliability problems of the
systems under evaluation. A second confound is contamination: when a
system states a correct law, standard evaluations cannot distinguish
derivation from memorization of published results.

We address both issues by construction. The benchmark plants laws that
are provable theorems, drawn from arithmetic statistics, so that
discovery can be scored exactly against a known answer key, and
generates task variants absent from the literature, so that
memorization and reasoning are separately identifiable.

\subsection*{Task family}

Let $K/\Q$ be a Galois number field with $G = \mathrm{Gal}(K/\Q)$, and
let $\zeta_K(s) = \sum_{n \ge 1} a_n(K)\, n^{-s}$ be its Dedekind zeta
function, so that $a_n(K)$ counts ideals of norm $n$. For a rational
prime $p$ with decomposition type $(e, f, g)$ in $K$ (ramification
index, inertia degree, and number of primes above $p$, with
$efg = [K{:}\Q]$), the $p$-Euler factor is $(1 - p^{-fs})^{-g}$, whence
\[
a_{p^k}(K) \;=\; \binom{g + k/f - 1}{\,g - 1\,} \text{ if } f \mid k,
\qquad a_{p^k}(K) = 0 \text{ otherwise},
\]
and $a_n$ is determined multiplicatively. The inertia degree $f$ at an
unramified prime is the order of its Frobenius class, which by
Chebotarev is equidistributed with respect to the counting measure on
$G$. Consequently the isomorphism type of $G$ leaves provable
statistical fingerprints in the coefficient vector
$(a_1, \dots, a_N)$: only element orders of $G$ occur as inertia
degrees, with prescribed densities. A representative criterion: for
$[K{:}\Q] = \ell^2$, the extension is cyclic iff $a_{p^\ell}(K) = 0$
for some prime $p$; only $C_{\ell^2}$ admits inert primes, and in the
noncyclic case $a_{p^\ell}$ is provably nonvanishing at every prime,
ramified or not \cite{LeeLee2025}. Each
benchmark task is a classification problem: given
$(a_1, \dots, a_N)$ for a field drawn from this generative model,
identify $G$ among the groups of the given order, and state the
distinguishing criterion. The setting follows Lee and
Lee~\cite{LeeLee2025}, who classified such $G$ with interpretable
models and promoted the learned decision logic to new theorems.

The benchmark distribution is the exact Chebotarev model: independent
Frobenius classes at each prime, i.e.\ the idealized unramified
statistics. This idealization is validated against genuine arithmetic:
on number fields drawn from the LMFDB (with zeta data computed from
defining polynomials via Dedekind's theorem), the proven quartic
criterion classifies at accuracy $1.000$, as the theorem guarantees.

\subsection*{Exact ground truth}

Because the generative model is known in closed form, each task ships
with quantities that are typically unobtainable:

\begin{itemize}
  \item \textbf{Bayes-optimal accuracy.} The class-conditional law of
  the observable prime signatures is an explicit categorical
  distribution per prime, so the Bayes risk is estimable to
  $\pm 5 \times 10^{-4}$ by direct signature sampling. Agent
  performance is reported as regret against this ceiling; the ceiling
  itself calibrates task difficulty (from $0.9998$ down to $0.84$
  across the suite).
  \item \textbf{Exact attribution ground truth.} The mutual information
  $I(a_n; G)$ is computable in closed form for every feature via the
  independence of prime-power factors. A system's stated feature
  attributions are scored against this ranking (rank correlation,
  precision at $k$), which exposes cases where predictive accuracy and
  explanation fidelity decouple; we observe decision trees within
  $0.3\%$ of the ceiling whose importances correlate weakly with the
  true ranking.
  \item \textbf{Reference theorems.} The correct explanation for each
  task is a proven statement (e.g.\ the inert-prime criterion above),
  serving as the answer key for conjecture grading.
\end{itemize}

A natural hardening step is formal verification of the answer keys.
The discrete content of each criterion (the element-order multisets of
the candidate groups, the vanishing pattern of the Euler coefficients
as a function of the inertia degree, and the resulting exclusion
arguments) is directly formalizable in Lean~4 against Mathlib, making
the ground truth machine-checked rather than merely peer-reviewed. The
analytic inputs (Euler products, the Chebotarev density theorem) would
enter as stated hypotheses pending their availability in Mathlib. We
view machine-checked answer keys as the natural endpoint of the
benchmark's premise: an evaluation is only as trustworthy as its
ground truth.

\subsection*{Evaluation protocols and contamination design}

Systems are evaluated under three protocols: \emph{rule induction}
(data and mathematical context provided; the system must return an
executable decision rule and a stated conjecture, with the rule scored
on held-out fields), \emph{derivation} (mathematical context only, no
data; tests whether the criterion can be derived rather than fit), and
\emph{direct classification} (few-shot in-context prediction). The
harness is model-agnostic: any prompt-to-text interface, including a
multi-agent system behind a single entry point, plugs in unchanged;
alternatively an agent may be given the dataset files and evaluated on
its returned artifacts. Sandboxed execution and oracle/majority
controls validate the pipeline independently of any model.

The generator requires only the element-order multiset of a candidate
group, so the suite includes tasks for degrees $14$, $21$, and $25$
whose distinguishing criteria follow from the same theory but, to our
knowledge, appear in no publication. Matched comparisons between
published and held-out tasks convert contamination from a hypothesis
into a measurement.

\subsection*{Preliminary findings}

A preliminary single-seed run of Claude Sonnet 4.5 exhibits the
separations the design targets. Under rule induction the model attains
near-ceiling regret on several tasks, including held-out degree $14$
(regret $0.007$). Under derivation it collapses to majority-class
performance on most tasks, with one exception: the degree-$9$ task,
whose criterion is published, where it achieves regret $0.004$. This
asymmetry is consistent with memorization of the published criterion;
the decisive comparison is against the structurally matched held-out
degree-$25$ task (both are the $\ell^2$ family), which is part of the
full multi-seed evaluation in progress. Notably, the model also scored
$0.958$ on the held-out degree-$14$ derivation, where no published
criterion exists; whether its stated conjecture constitutes a correct
derivation of the inert-prime argument is exactly what the
conjecture-grading protocol assesses, and the harness now records all
model outputs for this analysis.

\subsection*{Directions: a testbed for trustworthy multi-agent ML}

Beyond capability evaluation, the framework instruments security and
trust questions that are normally unmeasurable because harm cannot be
quantified without ground truth.

\emph{Poisoning with exact damage accounting.} Since the generative
model is known in closed form, training-distribution attacks can be
analyzed information-theoretically: perturbing the Chebotarev
distribution shifts the Bayes ceiling by a computable amount,
separating the information the attack destroys from the performance
the learner loses. Standard poisoning studies cannot make this
distinction.

\emph{Byzantine tolerance of agent pipelines.} In a multi-agent
discovery system, a compromised member can inject subtly flawed
analyses. With exact regret as the harm metric, degradation can be
plotted against the number and subtlety of compromised agents, and the
detection value of verifier roles measured directly. This is Byzantine
fault tolerance for epistemic outputs rather than messages.

\emph{Faithfulness auditing.} Gradeable explanations expose systems
that are right for the wrong reasons; we already observe models within
$0.3\%$ of the ceiling whose stated feature attributions correlate
weakly with the true information ranking. In a security setting this
is the difference between passing accuracy gates and being
trustworthy.

A first instrument for these questions is implemented: a tool-using
multi-agent pipeline (Strands agents on Amazon Bedrock) in which an
analyst agent explores the data and proposes a rule, and an
independent critic agent verifies or revises it. Four conditions
(analyst alone, analyst with critic, a simulated compromised analyst,
and the compromised analyst with critic) make the value of
verification and the harm of a compromised member directly comparable
on the same tasks; this evaluation is in progress. These questions
seem complementary to the APWA evaluation agenda.

Code, benchmark suite (18 configurations, Hugging Face format), and
harness: \repourl.

\begin{thebibliography}{1}

\bibitem{LeeLee2025}
K.-H. Lee and S. Lee,
\emph{Machines Learn Number Fields, But How? The Case of Galois Groups},
arXiv:2508.06670 [math.NT], 2025.
\url{https://arxiv.org/abs/2508.06670}

\end{thebibliography}

\end{document}
